%----------------------------------------------------------------------------
\chapter{\bevezetes}
%----------------------------------------------------------------------------

Az emberiség kommunikációjának védelme iránti igény egyidős a civilizációval. Az írott információ
titkosítása évezredek óta folyamatosan fejlődő diszciplína, amelyet minden korszak saját technológiai
szintjéhez igazított. A 20.\ század végén lezajlott digitális forradalom azonban minőségileg új
helyzetet teremtett: a kriptográfia többé nem pusztán katonai vagy diplomáciai eszköz, hanem a
globális digitális infrastruktúra alapköve. Bankrendszerek, egészségügyi nyilvántartások, ipari
vezérlőrendszerek és hétköznapi kommunikáció egyaránt a nyilvános kulcsú kriptográfia védelmére
támaszkodik. E rendszerek biztonsága matematikai problémák számítási nehézségére épül -- egy
feltevésre, amelynek érvényessége a közeljövőben komolyan megkérdőjeleződhet.

Párhuzamosan a digitális forradalom kibontakozásával egy másik technológiai változás is lezajlott:
a dolgok internete (Internet of Things, IoT) robbanásszerű terjedése. Beágyazott processzorok,
érzékelők és hálózati kapcsolat ma már nem csupán számítógépekben, hanem ipari gépekben, orvosi
implantátumokban, okos otthonokban és kritikus infrastruktúrában is megtalálható. A becslések
szerint 2030-ra több mint 25 milliárd hálózatba kötött eszköz lesz forgalomban~\cite{iot_forecast_2023},
amelyek mindegyike érzékeny adatokat generál, tárol és továbbít, miközben erőforrásaikban --
számítási teljesítményben, memóriában és energiában -- erősen korlátozottak.

A két tendencia találkozásánál egy harmadik, egyre fenyegetőbb fejlemény jelenik meg: a
kvantumszámítógépek fejlődése. A kvantummechanika elvein alapuló számítási paradigma alkalmas arra,
hogy a ma használt nyilvános kulcsú kriptográfiai rendszerek matematikai alapjait polinomiális időn
belül megoldja. Ez azt jelenti, hogy egy kellően erős kvantumszámítógép megjelenésével a jelenlegi
titkosítási infrastruktúra teljes egészében sérülékennyé válik -- nem csupán a jövőbeli
kommunikáció, hanem a ma titkosított és elmentett forgalom is, amelyet a jövőben visszafejtenek.

A dolgozat célja, hogy megvizsgálja, hogy hogyan lehet posztkvantum-kriptográfiai (PQC) megoldásokat --
konkrétan a CRYSTALS-Kyber és a BIKE kulcskapszulázó algoritmusokat -- hatékonyan alkalmazni
erőforrás-korlátozott IoT környezetekben. A bevezető fejezet az elméleti hátteret tárgyalja:
bemutatja a kriptográfia fejlődésének történetét, részletezi a kvantumszámítógépek működési elvét,
előnyeit és veszélyeit, majd ismerteti az IoT ökoszisztéma sajátosságait és biztonsági kihívásait,
végül bemutatja a PQC és IoT területek kapcsolatát.


%----------------------------------------------------------------------------
\section{A kriptográfia és az információbiztonság rövid története}
%----------------------------------------------------------------------------

A titkosírás legkorábbi ismert példái az ókori Egyiptomból származnak, ahol Kr.\ e.\ 1900 körül
módosított hieroglifákat alkalmaztak a szöveg értelmének elrejtésére. Az ókori Görögországban a
szkütalé nevű hengerrel transzpozíciós titkosítást végeztek, míg Julius Caesar a nevét viselő
Caesar-rejtjellel titkosította katonai üzeneteit: minden betűt az ábécé harmadik következő betűjével
helyettesített. Ezek az ún.\ klasszikus rejtjelek helyettesítésen vagy transzpozíción alapultak, azonban
a kriptoanalízis fejlődésével rendszeresen feltörhetővé váltak.

A középkorban a polialfabetikus rejtjelezés megjelenése (Vigenère-rejtjel, 1553) egy lépéssel
előrébb vitte a titkosítás nehézségét, ám a Kasiski-teszt (1863) és a Friedman-koincidensiindex
(1920-as évek) módszereivel ez a megközelítés is feltörhetőnek bizonyult. A klasszikus kriptográfia
lényegi korlátja az volt, hogy minden rendszer biztonsága a kulcs titkosságán alapult, miközben a
kulcsot valahogyan biztonságosan el kellett juttatni a kommunikáló felekhez -- ez az ún.\ kulcsterjesztési
probléma, amely évszázadokig megoldhatatlannak tűnt.

A modern kriptográfia alapjait a 20.\ század első felében fektették le. A második világháborúban az
Enigma-gép és az annak feltörésére irányuló bletchley-i erőfeszítések döntő szerepet játszottak a
konfliktus kimenetelében, és egyben megalapozták a számítástechnika és az algoritmikus kriptográfia
fejlődését. Alan Turing, Gordon Welchman és munkatársaik az elektromechanikus Bombe gépek
segítségével automatiulizálták az Enigma-kulcsok feltörését, és ezzel a kriptanalízist ipari méretűvé
tették. Claude Shannon 1949-es munkája, a \textit{Communication Theory of Secrecy Systems}
matematikai alapokra helyezte az információelmélet és a kriptográfia kapcsolatát: bevezette az
entrópia, a tökéletes titkosság és a biztonság formális fogalmait. Shannon
megmutatta, hogy tökéletes titkosság egyetlen rendszertől, az egyszer használatos rejtjeltől
(one-time pad) várható el, amely azonban kulcskezelési okokból nem praktikusan alkalmazható.

Az 1970-es évek hoztak igazi forradalmat. 1973-ban az USA Nemzeti Szabványügyi Intézete (NBS, ma
NIST) meghirdette az egységes adattitkosítási szabvány pályázatát, amelyből 1977-ben a DES (Data
Encryption Standard) lett. A DES egy 64 bites blokkméretű, 56 bites kulcsú szimmetrikus rejtjel,
amely az IBM Lucifer algoritmusán alapult, és amelynek tervezésébe az NSA is beleszólt. A DES
évtizedekig ipari szabvány maradt, ám 1998-ban az EFF (Electronic Frontier Foundation) \textit{Deep
Crack} nevű dedikált hardverével 56 óra alatt visszafejtett egy DES-titkosított üzenetet,
egyértelműen bizonyítva az algoritmus elavultságát. Utódja, a 2001-ben szabványosított AES (Advanced
Encryption Standard) -- Vincent Rijmen és Joan Daemen Rijndael algoritmusán alapulva -- máig az
egyik legelterjedtebb szimmetrikus titkosítási algoritmus.

Az igazi paradigmaváltást azonban a nyilvános kulcsú kriptográfia 1976-os bevezetése jelentette.
Whitfield Diffie és Martin Hellman \textit{New Directions in Cryptography} című munkájukban~\cite{diffie1976}
egy merőben új elvet fogalmaztak meg: a titkos kulcs megosztása nélkül is lehetséges biztonságos
kommunikáció létrehozása. Ez az ún.\ aszimmetrikus kriptográfia alapelve: egy nyilvános kulccsal
bárki titkosíthat, de visszafejteni csak a megfelelő privát kulcs birtokosa képes. A Diffie--Hellman
kulcscsere a diszkrét logaritmus problémán alapul, amelynek megoldása a csoport megfelelő
megválasztásával praktikusan kivitelezhetetlennek tűnt.

Ezt az ötletet rövidesen Ron Rivest, Adi Shamir és Leonard Adleman ültette át a gyakorlatba az RSA
algoritmus formájában (1977, közzétéve 1978-ban), amelynek biztonsága a nagy számok faktorizálásának
nehézségén alapul. Az 1980-as és 1990-es évtizedben Neal Koblitz és Victor Miller egymástól
függetlenül javasolta az elliptikus görbe kriptográfia (ECC) alkalmazását, ahol a csoportként
elliptikus görbék pontjait alkalmazzák. Az ECC kisebb kulcsmérettel nyújt azonos biztonsági szintet
az RSA-hoz képest, ami mobil és beágyazott alkalmazásokban komoly előnnyé vált.

A 2000-es évektől a kriptográfia az internet mindennapi infrastruktúrájává vált. Az új protokolloknak már alapkövetelménye a megfelelő szintű titkosítás. A titkosítást nem alkalmazó megoldásokat folyamatosan vezetik ki az internet forgalmából. Mára elképzeltetlennek tűnik egy olyan weboldal, amely nem használ HTTPS-t, vagy egy e-mail szolgáltató, amely nem támogatja a TLS-t. A VPN-ek, a biztonságos fájltárolás és a digitális aláírások mind a nyilvános kulcsú kriptográfia eredményei, amelyek nélkülözhetetlenek a modern digitális élethez. Azonban ez a globális infrastruktúra most szembesül azzal a fenyegetéssel, hogy a kvantumszámítógépek megjelenésével gyakorlatilag azonnal elavulttá válik.


%----------------------------------------------------------------------------
\section{Kvantumszámítógépek}
%----------------------------------------------------------------------------

A kvantumszámítógépek nem egyszerűen gyorsabb klasszikus számítógépek: egy alapvetően más számítási
paradigmát képviselnek, amely a kvantummechanika sajátos jelenségeit -- a szuperpozíciót, az
összefonódást és az interferenciát -- használja fel számítási erőforrásként. Ennek megértése
nélkülözhetetlen a kriptográfiai következmények értékeléséhez.

\subsection{A kvantumszámítás működési elvei}

A klasszikus számítástechnika alapegysége a bit, amely kizárólag a 0 vagy az 1 értéket veheti fel.
A kvantumszámítógép alapegysége a kvantumbit, röviden qubit, amelynek állapotát a kvantummechanika
szabályai írják le. A qubit állapotát egy kétdimenziós komplex Hilbert-tér elemeként ábrázolhatjuk legkönnyebben a Bloch-gömb segítségével. Egy qubit általános állapota egy szuperpozíció, amely a következőképpen írható fel:
$$\vert\psi\rangle = \alpha\vert 0\rangle + \beta\vert 1\rangle,$$
ahol $\alpha, \beta \in \mathbb{C}$ és $\vert\alpha\vert^2 + \vert\beta\vert^2 = 1$. Ez azt jelenti,
hogy mérés előtt a qubit egyszerre tartalmazza a $\vert 0\rangle$ és a $\vert 1\rangle$ állapotot,
az ún.\ szuperpozíció elvének megfelelően. Mérés hatására a szuperpozíció összeomlik: $\vert\alpha\vert^2$
valószínűséggel kapjuk a 0-t, $\vert\beta\vert^2$ valószínűséggel az 1-et.

A szuperpozíció önmagában még nem tenné hatékonnyá a kvantumszámítást. A másik kulcsfontosságú
jelenség az összefonódás (entanglement): két qubit összefonódott állapotban olyan módon korrelál,
hogy az egyik mérési eredménye azonnali hatással van a másikra, függetlenül a köztük lévő fizikai
távolságtól. Az összefonódás következménye, hogy $n$ qubit nem $2n$, hanem $2^n$ állapot egyidejű
reprezentációjára képes. Egy 300 qubites rendszer több lehetséges állapotot reprezentál egyszerre,
mint ahány atom van az ismert univerzumban.

A harmadik alapjelenség az interferencia. A kvantumalgoritmusokat úgy tervezik, hogy a helyes választ
adó amplitúdók konstruktívan interferáljanak (erősítsék egymást), a helytelen válaszokat adó
amplitúdók pedig destruktívan (gyengítsék egymást). Ez a mechanizmus teszi lehetővé, hogy a
kvantumalgoritmusok exponenciálisan nagy keresési tereken belül hatékonyan találják meg a megoldást.
A kvantumalgoritmus lényegében az amplitúdókat manipulálja kvantumkapukkal (unitér transzformációkkal),
majd a végállapot mérésével nyeri ki a hasznos információt.

A kvantumszámítógépek fizikai megvalósítása számos különböző technológiai úton halad. A szupravezető
qubitek (amelyeket az IBM, Google és mások fejlesztenek) mesterséges atomokat alkalmaznak
szupravezető Josephson-kontaktuson alapuló áramkörökben, kriogén hőmérsékleten ($\approx$15~mK).
A csapdázott ionok (IonQ, Quantinuum) elektromágneses csapdákban tartott ionokat használnak
qubitként, ami hosszabb koherenciaidőt és alacsonyabb hibarátát eredményez. A fotonikus és
topologikus kvantumszámítógépek szintén kutatás alatt állnak. 2019-ben a Google bejelentette a
kvantumfölény elérését: a Sycamore processzor 53 qubittel egy mintavételi feladatot 200 másodperc
alatt oldott meg, amelyhez a legjobb klasszikus szuperszámítógépek becsléseik szerint 10\,000 évet
igényeltek volna~\cite{google_supremacy_2019}. 2023-ra az IBM Eagle és Heron processzorai már 127,
illetve 133 qubitet tartalmaztak, és a hibajavítással ellátott logikai qubitek is megjelentek.

A kvantumhiba-javítás a skálázhatóság egyik legnagyobb kihívása. Fizikai qubitek hibával terheltek:
a koherenciaidő véges, a kapuhiba valószínűsége $10^{-3}$--$10^{-4}$ nagyságrendű. A hibajavított
logikai qubitek eléréséhez becslések szerint 1\,000--10\,000 fizikai qubit szükséges egyetlen
logikai qubitenként, ami azt jelenti, hogy egy kriptográfiailag releváns, néhány ezer logikai
qubitet igénylő számítás elvégzéséhez millió fizikai qubit méretű gép szükséges. Ez ma még nem
áll rendelkezésre, de a fejlesztési ütem gyors.

\subsection{A kvantumszámítógépek előnyei és lehetőségei}

A kvantumszámítástechnika potenciális alkalmazásai messze túlmutatnak a kriptográfiai fenyegetésen,
és számos területen ígérik a klasszikus számítástechnika fundamentális korlátainak áttörését.

A \textbf{molekulaszimulációk és kvantumkémia} területén a kvantumszámítógépek képesek
molekulák elektronszerkezetét kvantummechanikai szinten szimulálni. Míg klasszikus számítógépen
egy közepes méretű molekula pontos szimulációja exponenciálisan növekvő erőforrást igényel,
addig egy kvantumszámítógép természetes hatókörébe esik a feladat. Ez közvetlen alkalmazást ígér
a gyógyszerfejlesztésben (hatóanyag--fehérje kölcsönhatások), az anyagtudományban (új szupravezetők,
akkumulátor anyagok tervezése) és a katalízisvizsgálatban.

\textbf{Optimalizálási feladatok} esetén a kvantum heuristikák -- mint a Quantum Approximate
Optimization Algorithm (QAOA) és a Quantum Annealing -- NP-nehéz problémák közelítő megoldására
mutatnak ígéretes eredményeket. Logisztika, ellátásilánc-tervezés, pénzügyi portfólió-optimalizálás
és menetrendtervezés területén már kisebb méretű kísérletekben is demonstráltak kvantum
előnyt.

A \textbf{pénzügyi modellezésben} a Monte Carlo-szimulációk kvantum gyorsítása négyzetgyök-csökkentést
ígér a mintaszámban, ami kockázatbecslés és derivatív árazás területén jelent versenyelőnyt.
Az \textbf{mesterséges intelligencia} területén kvantum-neuronhálózatokat és kvantumos kernel-módszereket
vizsgálnak, bár itt az előny kevésbé egyértelmű, mint az előző területeken.

Végül meg kell említeni a \textbf{kvantumkulcs-elosztást (QKD)}, amely a kvantummechanika
törvényein alapuló biztonságos kulcscserét kínál: a kvantumállapotok lemásolhatatlansága
(no-cloning tétel) és a mérés rendelkező hatása garantálja, hogy a lehallgatás kimutatható.
A QKD azonban saját infrastruktúrát és speciális hardvert igényel, és nem közvetlenül helyettesíti
a PQC-t.

\subsection{Veszélyek a jelenlegi kriptográfiai rendszerekre}

A kvantumszámítógépek kriptográfiai fenyegetése két alapvető algoritmuson keresztül valósul meg.

\textbf{A Shor-algoritmus:} 1994-ben Peter Shor bebizonyította, hogy egy kellően erős kvantumszámítógép polinomiális
idő alatt képes az egészfaktorizáció és a diszkrét logaritmus problémájának megoldásá\-ra~\cite{shor1994}.
Az algoritmus kvantum-Fourier-transzformációt alkalmaz az adott szám prímtényezőinek
meghatározásához. A faktorizálás időbonyolultsága $O((\log N)^3)$, szemben a legjobb ismert
klasszikus szub-exponenciális algoritmussal, az általános számmező szitával (GNFS), amely körülbelül
$O(\exp((\log N)^{1/3}))$ nagyságrendű. Ez közvetlenül aláássa az RSA, a Diffie--Hellman
kulcscsere és az ECC biztonságát: az RSA-2048 feltöréséhez hozzávetőleg 4\,000 hibajavított
logikai qubit szükséges -- ez a szám a jelenlegi fejlesztési ütem alapján akár a 2030-as évekre
el is érhető.

\textbf{A Grover-algoritmus:} 1996-ban Lov Grover megmutatta, hogy egy $N$ elemű rendezetlen
adatbázisban a keresés kvantumszámítógépen $O(\sqrt{N})$ lépésben elvégezhető,
szemben a klasszikus $O(N)$ lépéssel. Szimmetrikus kriptográfiánál ez a négyzetgyök-gyorsulás azt
jelenti, hogy egy 128 bites AES kulcs bruteforce-keresése olyan, mintha 64 bites kulcsra kellene
keresni: a megoldás a kulcsméretek megduplázása, azaz az AES-256 alkalmazása. A szimmetrikus
kriptográfia tehát módosítással kvantum-ellenálló maradhat.

A legjelentősebb stratégiai veszélyt azonban a \textbf{„gyűjtsd most, fejtsd meg később"}
(harvest-now, decrypt-later, HNDL) stratégia jelenti. Állami szintű szereplők és fejlett
kiberbűnözői csoportok már most gyűjthetik a titkosított hálózati forgalmat, hogy azt egy
jövőbeli kvantumszámítógéppel visszafejtsék. Azok az adatok, amelyeknek ma titkosítva kell
maradniuk 10--20 éves időtávon (állami titkok, egészségügyi adatok, szellemi tulajdon,
pénzügyi tranzakciók), már most célponttá válnak. Ez a fenyegetés sürgetővé teszi a PQC
bevezetését még azelőtt, hogy a kvantumszámítógépek kriptográfiai szintre érnek.

A NIST 2016-ban meghirdetett versenye ennek a felismerésnek a közvetlen következménye volt.
A hat éves értékelési folyamat eredményeként 2022-ben a NIST közzétette az első négy
posztkvantum-kriptográfiai algoritmust: a CRYSTALS-Kyber-t (kulcskapszulázás), a
CRYSTALS-Dilithium-ot és a FALCON-t (digitális aláírás), valamint a SPHINCS+-t (hash-alapú
aláírás)~\cite{NIST-Select-2022}. 2024-ben ezek FIPS szabványként is megjelentek (FIPS 203, 204, 205).


%----------------------------------------------------------------------------
\section{A dolgok internete}
%----------------------------------------------------------------------------

Az IoT fogalma egy mélyen heterogén ökoszisztémát takar, amelynek határai az évek során
folyamatosan tágulnak. Az egyszerű otthoni okoseszközöktől az ipari automatizálás
vezérlőrendszerein át az orvosi implantátumokig terjed az eszközök köre, amelyek összekötik a
fizikai és a digitális világot.

\subsection{Az IoT kialakulása és fejlődése}

Az ,,összes dolgot hálózatba kötni'' ötlete valójában megelőzi az ,,Internet of Things'' elnevezést.
Mark Weiser, a Xerox PARC kutatója 1991-es \textit{The Computer for the 21st Century} esszéjében
lefektette az ubiquitous computing (mindenütt jelenlévő számítástechnika) koncepcióját: elképzelte,
hogy számítástechnikai eszközök a mindennapi tárgyakba integrálódnak, és a felhasználó számára
láthatatlanul végzik munkájukat. Az 1990-es évek elején a Carnegie Mellon
Egyetemen már létezett az első hálózatba kötött automatával bemutatott kísérlet, amely szenzorok
segítségével tájékoztatta a felhasználókat az üdítőgép aktuális tartalmáról.

A ,,dolgok internete'' elnevezést Kevin Ashton brit technológus alkotta meg 1999-ben, eredetileg
az RFID (rádiófrekvenciás azonosítás) és az ellátásilánc-menedzsment összefüggésében. Az elgondolás
lényege, hogy ha minden fizikai tárgyhoz egyedi azonosítót és hálózati kapcsolatot rendelünk, az
objektumok maguk is adatforrásokká válnak, és a valós világ állapota automatikusan nyomon
követhetővé válik. Az RFID technológia lehetővé tette például a raktári
készletezés automatizálását, amelyet a Procter \& Gamble és az Auto-ID Center együttműködésével
fejlesztett ki.

A 2000-es évek elejétől az ipari szenzorhálózatok (Wireless Sensor Networks, WSN) terjedtek el a
gyártásban, a mezőgazdaságban és a katonai megfigyelésben. A fordulópontot az okostelefonok
megjelenése (2007), az IPv6 szabvány széleskörű bevezetése, valamint az alacsony fogyasztású
rádióprotokollok (Bluetooth LE, ZigBee, LoRaWAN, NB-IoT) és az olcsó ARM-alapú mikrovezérlők
tömeges piaci megjelenése hozta el. 2010-től az IoT eszközök száma exponenciálisan növekedett:
2015-re elérte a 15 milliárdot, 2020-ra meghaladta a 30 milliárdot.

Az intelligens otthonok piacán a Nest termosztát (2011), az Amazon Echo (2014) és a Philips Hue
izzósorozat paradigmaváltást hozott a fogyasztói elektronikában. Az ipari IoT a gyártásban és az
energetikában a prediktív karbantartás, a folyamatoptimalizálás és a valós idejű monitorozás
eszközévé vált. Az egészségügyi IoT területén a viselhető szenzoros eszközöktől a beültetett
pacemaker-eken és szívritmus-figyelőkön át az infúziós pumpák hálózati vezérléséig terjedt az
alkalmazás. A McKinsey Global Institute becslése szerint az IoT gazdasági hatása 2030-ra
5,5--12,6 billió dollárt tesz ki évente~\cite{mckinsey_iot}.

\subsection{IoT eszközök jellemzői és korlátai}

Az IoT ökoszisztéma egyik legfontosabb sajátossága a rendkívüli heterogenitás és az erőforrás-korlátok
jelenléte. Egy nagyvállalati szerver és egy ipari hőmérséklet-érzékelő között több nagyságrendnyi
különbség lehet számítási teljesítményben, memóriában és energiaellátásban.

Tipikus korlátok az alacsony kategóriás IoT eszközöknél:
\begin{itemize}
    \item \textbf{Processzor:} 8--32 bites mikrovezérlők (pl.\ ARM Cortex-M0/M4, Atmel AVR),
          órajel 8--120~MHz tartományban, hardveres lebegőpontos egység nélkül.
    \item \textbf{Memória:} 2--256~kB RAM, néhány kB -- néhány MB flash tároló.
    \item \textbf{Energia:} elemről működő eszközöknél az aktív kriptográfiai számítás az egyik
          legnagyobb energiaigényű művelet, ezért fontos a kulcscsere-műveletek ritkítása.
    \item \textbf{Sávszélesség:} sok protokoll (LoRaWAN, NB-IoT) esetén a payload mérete
          néhány tíz--száz bájtra korlátozott, ami a PQC kulcs- és kapszulaméreteit
          érzékenyen érinti.
    \item \textbf{Üzemidő:} ipari és infrastrukturális eszközöket 10--20 éves élettartamra
          terveznek, ami azt jelenti, hogy a ma telepített eszközök a kvantumfenyegetés
          kritikus időszakában is aktívan üzemelni fognak.
\end{itemize}

Ez utóbbi szempont különösen fontos. Egy 2024-ben telepített, 15 éves élettartamra tervezett
ipari szenzor 2039-ig működik. Ha a 2030-as évekre valóban megjelennek kriptográfiailag releváns
kvantumszámítógépek, ez az eszköz az aktív üzemeltetési ideje alatt válik sebezhető célponttá.
Az ún.\ \textit{cryptographic agility} -- a kriptográfiai algoritmusok futás közbeni cserélhetőségének
megőrzése a rendszerarchitektúrában -- ezért az IoT tervezés egyik alapkövetelményévé vált.

\subsection{IoT biztonsági kihívások}

Az IoT biztonsági kérdései már jóval a kvantumfenyegetés előtt komoly figyelmet kaptak. A 2016-os
Mirai botnet-támadás megmutatta, hogy hétköznapi IoT eszközök (IP-kamerák, DVR-ek) gyári
jelszóval hagyott eszközei milyen pusztító erejű DDoS-támadáshoz nyújthatnak alapot: a
botnet csúcsán másodpercenként 620~Gbps forgalmat generált, és az internetes infrastruktúra
jelentős részét megbénította~\cite{mirai2016}. A 2021-es Oldsmar vízkezelő-üzemi incidens, ahol
egy támadó a klórszint módosításával próbálkozott, jól illusztrálta az IoT rendszerek fizikai
következményekkel járó sebezhetőségét.

A klasszikus kriptográfiai fenyegetéseken túl az IoT specifikus biztonsági kihívásai a következők:

\textbf{Hitelesítési problémák.} Erőforrás-korlátozott eszközökön a PKI alapú tanúsítványkezelés
(TLS 1.3 futtatása, tanúsítványmegújítás, visszavonás-ellenőrzés) számottevő CPU-terhelést és
memóriát igényel. Sok eszköz ezért gyári, nehezen cserélhető, megosztott titkokra támaszkodik,
ami egypontos sebezhetőséget teremt.

\textbf{Fizikai hozzáférés és oldalsáv-támadások.} Az IoT eszközök fizikailag hozzáférhető helyen
üzemelnek, fizikai biztonságuk korlátozott. A fogyasztásmérés-alapú (Differential Power Analysis,
DPA) és elektromágneses sugárzáson alapuló (Electromagnetic Analysis, EMA) elemzések titkok
kinyerésére alkalmasak, hacsak az implementáció nem állandó idejű (constant-time) és nem alkalmaz
megfelelő zaj-ellenintézkedéseket.

\textbf{Firmware-frissítések és cryptographic agility.} A hosszú üzemidő miatt az algoritmusváltásnak
szoftverfrissítésen keresztül kell megvalósulnia. Ez előre megtervezett biztonságos bootloadereket,
aláírt firmware-csomagokat és megbízható frissítési csatornákat feltételez. Ezek hiányában a
deployed eszközök nem frissíthetők kvantum-ellenálló algoritmusokra.

\textbf{Kulcskezelés és üzenetméret.} Az IoT rendszerek skálája -- sok milliós eszközpark --
egyedi kulcsmenedzsment kihívásokat vet fel. A PQC implementáció során figyelemmel kell lenni
arra, hogy a rács alapú és kód alapú algoritmusok nyilvános kulcsai és kapszulái lényegesen
nagyobbak, mint az RSA vagy ECC megfelelőik, ami ütközhet az eszközök korlátozott tárhelyével
és a sávszélességkorlátokkal.


%----------------------------------------------------------------------------
\section{Posztkvantum-kriptográfia és IoT -- a kapcsolat}
%----------------------------------------------------------------------------

A posztkvantum-kriptográfia azokat a kriptográfiai algoritmusokat foglalja magában, amelyek
ellenállnak mind a klasszikus, mind a kvantumszámítógépek által nyújtott támadásoknak. Elnevezése
némileg félrevezető: a ,,poszt-'' előtag nem azt jelenti, hogy a kvantumszámítógépek már megjelentek,
hanem azt, hogy az algoritmusok a kvantumos utáni korszakra is biztonságosak.

A PQC főbb matematikai megközelítései:
\begin{itemize}
    \item \textbf{Rács alapú kriptográfia:} az SVP és az LWE probléma nehézségén alapul.
          Ide tartoznak a CRYSTALS-Kyber (KEM) és CRYSTALS-Dilithium (aláírás),
          amelyeket a NIST 2022-ben szabványosított~\cite{NIST-Select-2022}.
    \item \textbf{Kód alapú kriptográfia:} lineáris hibakijavító kódok dekódolásának
          NP-nehézségén alapul. A BIKE és a HQC ide sorolható; a McEliece-rendszer
          1978 óta ismert, de kulcsméretei igen nagyok.
    \item \textbf{Hash alapú kriptográfia:} kiforrottsága és konzervatív biztonsági
          feltevései miatt kisebb aláírási rendszerek esetén vonzó (SPHINCS+).
    \item \textbf{Izogénián alapuló kriptográfia:} a SIKE algoritmust 2022-ben klasszikus
          támadással feltörték, ami e kategória korlátjaira hívja fel a figyelmet.
\end{itemize}

Az IoT kontextusában a PQC bevezetése különleges kihívásokat jelent. A CRYSTALS-Kyber-512
esetén a nyilvános kulcs 800 bájt, a kapszula 768 bájt, míg az RSA-2048 esetén ez 256 bájt --
a háromszoros méretbeli növekedés mind a sávszélességet, mind a flash-tárhelyt igénybe veszi.
Az elterjedt TLS és DTLS protokollok üzenetméreteinek módosítására is szükség van.

A számítási igény tekintetében kedvezőbb a kép: a CRYSTALS-Kyber és a BIKE ARM Cortex-M4
processzorral néhány tíz--száz milliszekundum alatt elvégezhetők~\cite{CRYSTALS-KYBER, BIKE-SPEC},
ami a legtöbb IoT alkalmazásban elfogadható. A \textbf{hibrid megközelítés} -- klasszikus és
PQC kulcscsere párhuzamos alkalmazása -- a biztonsági migráció kockázatát is csökkenti: a
kommunikáció biztonságos marad mindaddig, amíg legalább az egyik algoritmus ellenáll
a fenyegetésnek. Ez ma az iparági ajánlás a PQC fokozatos bevezetésekor.

Az ipari és szabályozói nyomás is egyre erősebb. Az Európai Unió NIS2-irányelve és a Cyber
Resilience Act kötelezővé teszik a beágyazott rendszerek gyártói számára a biztonsági
frissíthetőség biztosítását és a kvantum-ellenálló kriptográfiára való tervezett átállást.
Az IETF aktívan dolgozik a PQC-t integráló TLS és DTLS profilok kidolgozásán, különösen a
CoAP és a 6LoWPAN kombinálásával működő IoT-profilok esetén~\cite{NIST-PQC}.

Az összefonódó kihívások -- a kvantumfenyegetés, az IoT korlátok és a szabályozói nyomás --
együttesen sürgetik a PQC algoritmusok erőforrás-tudatos implementációinak fejlesztését. E dolgozat
éppen erre a hiányra ad választ: a CRYSTALS-Kyber és a BIKE algoritmusok IoT kontextusban való
teljesítményének részletes mérésével és elemzésével.


%----------------------------------------------------------------------------
\section{A dolgozat céljai és felépítése}
%----------------------------------------------------------------------------

A dolgozat fő kutatási kérdése: \textit{A CRYSTALS-Kyber és BIKE posztkvantum-kriptográfiai
algoritmusok erőforrás-korlátozott IoT eszközökön való alkalmazhatóságának és egymáshoz
viszonyított teljesítménybeli különbségeinek meghatározása szimulált mérési körülmények között.}

A kutatási célkitűzés megvalósításához a következő részfeladatokat tűztem ki:
\begin{enumerate}
    \item A CRYSTALS-Kyber és a BIKE algoritmusok elméleti hátterének és matematikai
          alapjainak részletes ismertetése, különös tekintettel az IoT-szempontból
          releváns implementációs kihívásokra.
    \item Mérési környezet kialakítása, amely lehetővé teszi a két algoritmus
          szimulált erőforrás-korlátozott eszközön való teljesítmény-összehasonlítását.
    \item A kulcsgenerálás, a kulcskapszulázás (encapsulation) és a visszafejtés
          (decapsulation) futási idejének, memóriaigényének és energiaszükségletének
          mérése és kiértékelése.
    \item Következtetések levonása arról, melyik algoritmus alkalmasabb egyes IoT
          alkalmazási területekre, és milyen paraméterbeállítások optimálisak a
          biztonság és az erőforrásigény közötti kompromisszum szempontjából.
\end{enumerate}

